\section*{§1. The Field of Complex Numbers}

\textbf{Definition 1.} A \textit{complex number} is an expression of the form
\[
z = x + iy,
\]
where $x$ and $y$ are real numbers, and $i$ is a symbol satisfying the relation
\[
i^{2} = -1.
\]
The real numbers $x$ and $y$ are called the \textit{real part} and the \textit{imaginary part} of the complex number $z$, respectively. We write
\[
x = \text{Re } z, \quad y = \text{Im } z.
\]
Two complex numbers $z_{1} = x_{1} + iy_{1}$ and $z_{2} = x_{2} + iy_{2}$ are said to be \textit{equal} if and only if $x_{1} = x_{2}$ and $y_{1} = y_{2}$.

The set of all complex numbers is denoted by $\mathbb{C}$.

\textbf{Definition 2.} The \textit{sum} and the \textit{product} of two complex numbers $z_{1} = x_{1} + iy_{1}$ and $z_{2} = x_{2} + iy_{2}$ are defined by the formulas
\begin{align*}
z_{1} + z_{2} &= (x_{1} + x_{2}) + i(y_{1} + y_{2}), \\
z_{1} z_{2} &= (x_{1}x_{2} - y_{1}y_{2}) + i(x_{1}y_{2} + x_{2}y_{1}).
\end{align*}

It is easy to check that the operations of addition and multiplication of complex numbers satisfy the following properties:
\begin{enumerate}
\item[(i)] $z_{1} + z_{2} = z_{2} + z_{1}$, $z_{1}z_{2} = z_{2}z_{1}$ \hfill (commutativity);
\item[(ii)] $(z_{1} + z_{2}) + z_{3} = z_{1} + (z_{2} + z_{3})$, $(z_{1}z_{2})z_{3} = z_{1}(z_{2}z_{3})$ \hfill (associativity);
\item[(iii)] $z_{1}(z_{2} + z_{3}) = z_{1}z_{2} + z_{1}z_{3}$ \hfill (distributivity).
\end{enumerate}
The complex number $0 = 0 + i0$ is the \textit{zero element} of $\mathbb{C}$, since
\[
z + 0 = z \quad \text{for all } z \in \mathbb{C}.
\]
The complex number $1 = 1 + i0$ is the \textit{unit element} of $\mathbb{C}$, since
\[
z \cdot 1 = z \quad \text{for all } z \in \mathbb{C}.
\]
For every complex number $z = x + iy$, the complex number $-z = -x - iy$ is the \textit{opposite} of $z$, since
\[
z + (-z) = 0.
\]